\chapter{RSA}

\chapterSubhead{The padlock that secures the financial internet---and the one quantum will pick}

In 1977 Rivest, Shamir, and Adleman published the cryptographic protocol now universally called \keyterm{RSA} \citep{rivest_method_1978}. Half a century later, RSA is the most widely deployed public-key cryptosystem on earth: TLS handshakes, code signing, smartcards, S/MIME email, IPsec VPNs, financial settlement protocols, government identity documents. Almost every encrypted message you have ever sent has, at some point in its life, depended on RSA. This chapter explains how RSA works, why it has held up so well against classical cryptanalysis, and why Shor's algorithm (Chapter on Shor) is its decisive enemy.

%------------------------------------------------------------------
\section{Public-key cryptography in one paragraph}
%------------------------------------------------------------------

Until 1976, all known cryptosystems required sender and receiver to share a secret key in advance. \citet{diffie_new_1976} proposed the radical alternative: each party has a \emph{public} key that anyone can use to encrypt, and a separate \emph{private} key that only the owner can use to decrypt. Encryption is easy; decryption without the private key is computationally hopeless. This is \keyterm{public-key} or \keyterm{asymmetric} cryptography. Diffie--Hellman provided the first such protocol (for key exchange); RSA \citep{rivest_method_1978}, a year later, provided one for both encryption and digital signatures.

The breakthrough rested on the idea of a \keyterm{trapdoor function}: easy to compute in one direction, hard to invert without secret information. RSA's trapdoor is modular exponentiation: $x \mapsto x^e \bmod N$ is fast; $y \mapsto y^{1/e} \bmod N$ is, without the private key, equivalent in difficulty to factoring $N$.

%------------------------------------------------------------------
\section{The RSA construction}
%------------------------------------------------------------------

\textbf{Key generation.}
\begin{enumerate}
\item Choose two large primes $p$ and $q$. Compute $N = pq$. $N$ becomes part of the public key.
\item Compute Euler's totient $\varphi(N) = (p-1)(q-1)$.
\item Choose a public exponent $e$ with $\gcd(e,\varphi(N)) = 1$. Common choices: $e = 65537 = 2^{16} + 1$ (efficient for verification).
\item Compute the private exponent $d \equiv e^{-1}\pmod{\varphi(N)}$ via the extended Euclidean algorithm.
\item Public key: $(N, e)$. Private key: $d$ (and possibly $p, q$).
\end{enumerate}

\textbf{Encryption.} For a message $m$ with $0 \leq m < N$, compute ciphertext $c = m^e \bmod N$.

\textbf{Decryption.} Compute $m = c^d \bmod N$.

The correctness of decryption rests on Fermat's little theorem (generalised via Euler): for any $m$ coprime to $N$, $m^{\varphi(N)} \equiv 1 \pmod{N}$. Since $ed \equiv 1 \pmod{\varphi(N)}$, we have $ed = 1 + k\varphi(N)$ for some integer $k$, so
\[
  c^d = (m^e)^d = m^{ed} = m^{1+k\varphi(N)} = m\cdot(m^{\varphi(N)})^k \equiv m \pmod{N}.
\]
The encryption and decryption operations cancel.

\begin{example}
Pick $p=61$, $q=53$. Then $N = 3233$, $\varphi(N) = 60\cdot 52 = 3120$. Choose $e = 17$; then $d = e^{-1} \bmod 3120 = 2753$. To encrypt $m = 65$: $c = 65^{17} \bmod 3233 = 2790$. To decrypt: $m = 2790^{2753} \bmod 3233 = 65$. Tiny example; real RSA uses primes hundreds of digits long.
\end{example}

%------------------------------------------------------------------
\section{Where the security comes from}
%------------------------------------------------------------------

RSA's security rests on the difficulty of \keyterm{integer factorization}: given $N = pq$ with $p,q$ unknown large primes, recover $p$ and $q$. If an adversary could factor $N$, they could compute $\varphi(N) = (p-1)(q-1)$, then $d = e^{-1} \bmod \varphi(N)$, and decrypt at will.

The reduction goes only one way (factoring is sufficient to break RSA). It is an open question whether breaking RSA \emph{requires} factoring, or whether a separate algorithm could break RSA without recovering $p$ and $q$. For practical purposes, factoring is the assumed hardness.

Factoring is hard \emph{classically}. The best general-purpose algorithm is the \keyterm{number field sieve} (NFS) \citep{lenstra_factoring_1993}, with sub-exponential complexity
\[
  L_N\!\left[\tfrac{1}{3}, \sqrt[3]{\tfrac{64}{9}}\right] = \exp\!\Big((c+o(1))\,(\log N)^{1/3}(\log\log N)^{2/3}\Big), \quad c=\sqrt[3]{64/9}\approx 1.923.
\]
This is faster than fully exponential time but much slower than polynomial. For a 2048-bit RSA key, NFS requires somewhere in the neighbourhood of $10^{30}$ operations; the largest factored RSA modulus to date is RSA-250 (829 bits), taking 2700 core-years of compute in 2020. A 2048-bit modulus is, by every classical projection, secure for decades.

\begin{remark}
Note ``classically secure for decades.'' The same modulus on a sufficiently large quantum computer running Shor's algorithm is broken in hours. The asymmetry between classical sub-exponential and quantum polynomial complexity is what makes the post-quantum cryptography problem urgent.
\end{remark}

%------------------------------------------------------------------
\section{RSA in the financial system}
%------------------------------------------------------------------

A non-exhaustive list of the places RSA secures the financial infrastructure:

\begin{itemize}
  \item \textbf{TLS.} Web browsers and banking apps use RSA in the TLS handshake to establish symmetric session keys. Every online banking session begins with an RSA-protected key exchange.
  \item \textbf{Card networks.} EMV chip cards use RSA-based digital signatures to authenticate themselves to point-of-sale terminals; Visa and Mastercard issuer authentication relies on RSA.
  \item \textbf{SWIFT.} The interbank messaging network uses RSA-based public-key infrastructure for message authentication and key distribution.
  \item \textbf{Digital signatures on contracts.} S/MIME and PKCS\#7 signatures on legal documents, securities filings, regulatory submissions almost always use RSA.
  \item \textbf{Code signing.} Software updates to bank applications and trading systems are signed with RSA to prevent tampering.
\end{itemize}

The pervasiveness is the problem. Replacing RSA throughout this infrastructure requires coordinated migration across thousands of organizations, decades of legacy code, embedded hardware with fixed cryptography, and protocols that pre-date the post-quantum cryptography standards. \citet{auer_quantum_2024} estimate that a financial-sector-wide migration will take 10 years or more even with concerted effort.

%------------------------------------------------------------------
\section{Variants and accelerators}
%------------------------------------------------------------------

Several variants and accelerators have been developed for RSA over the years.

\textbf{RSA-OAEP.} Plain ``textbook'' RSA is malleable and deterministic, and is therefore not used in production. Real RSA encryption uses Optimal Asymmetric Encryption Padding (OAEP), which adds randomness and structure to make the ciphertext indistinguishable from random.

\textbf{RSA-PSS.} The probabilistic signature scheme adds a random salt to RSA signatures, providing stronger security proofs.

\textbf{CRT decryption.} Decryption can be sped up by a factor of $\sim 4$ using the Chinese Remainder Theorem to compute $c^d \bmod p$ and $c^d \bmod q$ separately and combine.

\textbf{Multi-prime RSA.} Using three or more primes can speed decryption further at a modest security cost.

None of these variants change the fundamental quantum vulnerability: all rely on the hardness of factoring or related number-theoretic problems that quantum computers can solve in polynomial time.

%------------------------------------------------------------------
\section{Pollard rho and other classical factoring techniques}
%------------------------------------------------------------------

For completeness, a few classical factoring techniques worth knowing:

\textbf{Trial division.} Test small primes as factors. Works for $N$ with small prime factors; fails immediately for products of two large primes (the RSA case).

\textbf{Pollard rho} \citep{pollard_monte_1975}. A clever random-walk approach that runs in $O(N^{1/4})$ time. Practical for moderately sized $N$ but quickly outpaced by the number field sieve.

\textbf{Pollard $p-1$.} Effective when $p-1$ has only small prime factors; mitigated by choosing $p$ such that $p-1$ has at least one large prime factor.

\textbf{Quadratic sieve.} Faster than Pollard for $N$ above $\sim 10^{50}$, slower than the number field sieve above $\sim 10^{100}$.

\textbf{Number field sieve.} The current champion for large $N$, the algorithm that has factored the largest RSA challenge moduli to date.

The classical state of the art improves slowly. Each algorithmic advance shifts the secure key size somewhat, but the asymptotic complexity has not changed since NFS in the early 1990s. By contrast, quantum factoring (Shor) is asymptotically polynomial---a category change, not an incremental improvement.

%------------------------------------------------------------------
\section{The harvest-now--decrypt-later threat}
%------------------------------------------------------------------

A practical concern often overlooked: an adversary does not need a quantum computer \emph{today} to attack RSA-encrypted data \emph{today}. They need only to record the encrypted traffic now and decrypt it later when a quantum computer becomes available. This is the \keyterm{harvest now, decrypt later} threat.

For information that loses value quickly---one-time passwords, ephemeral chat---this is not a concern. For information that retains value for years or decades---strategic positions, intellectual property, classified intelligence, long-term financial records, identity-document signatures---it is an active threat right now. State-level adversaries are widely believed to be archiving encrypted internet traffic precisely for this purpose.

\citet{auer_quantum_2024} highlight this point for the financial sector: client information, contract clauses, regulatory filings, and trade secrets routinely require multi-year or multi-decade confidentiality. Any such data encrypted with RSA today is at risk if quantum computers become powerful enough during its required confidentiality window.

%------------------------------------------------------------------
\section{The post-quantum migration}
%------------------------------------------------------------------

The cryptographic community has spent fifteen years developing public-key schemes that resist both classical and quantum cryptanalysis: \keyterm{post-quantum cryptography} (PQC). NIST ran a multi-round PQC competition from 2016 to 2024, culminating in the standardisation of:

\begin{itemize}
  \item \textbf{ML-KEM} (Module-Lattice Key Encapsulation Mechanism), based on the Module-LWE problem \citep{nist_fips203_2024}. Replaces RSA for key exchange.
  \item \textbf{ML-DSA} (Module-Lattice Digital Signature Algorithm). Replaces RSA for digital signatures.
  \item \textbf{SLH-DSA} (Stateless Hash-Based Digital Signature Algorithm). A hash-based signature alternative.
\end{itemize}

These are now FIPS standards (FIPS 203, 204, 205). Major browser vendors, cloud providers, and financial standards bodies are rolling them out. Financial institutions are advised to follow NIST's roadmap: inventory cryptographic dependencies, prioritize high-value long-confidentiality data, deploy hybrid (classical + post-quantum) protocols during transition, fully retire RSA in production over the next 5--10 years.

The migration is large and complex but tractable. The chapters on \keyterm{elliptic curves} and \keyterm{Shor's algorithm} complete this thread: ECC, the other major public-key family, faces the same fate as RSA, and Shor's algorithm is the single quantum reason why.

\textsep
